\documentclass[aps,prl,twocolumn,10pt,groupedaddress,nopacs]{revtex4}
\usepackage{amsmath,amssymb,graphicx}

\begin{document}

\title{The Geometric Standard Model: Unification via Reflexive Lattice Dynamics}
\author{William J. Steinmetz III}
\affiliation{Independent Researcher}
\date{\today}

\begin{abstract}
We present a unified field theory where the fundamental constants and mass spectrum of the Standard Model emerge from the geometric constraints of a discrete, reflexive spacetime lattice. We introduce a Reflexive Lagrangian with temporal non-locality, \textbf{proving} that the vibrational spectrum is quantized in powers of the Golden Ratio. This dynamical mechanism \textbf{generates} the integer partition $\{7, 3, 13, 4\}$ identified as the ``Integer Bootstrap.'' Using this framework, we derive the fine-structure constant $\alpha^{-1} = 137.036$ (1.3 ppm accuracy) from the 16 degrees of freedom of the minimal cubic cell. We extend the baryon mass formula to excited states, predicting the Delta Baryon mass at 1231.7 MeV (0.03\% error) without free parameters. The Higgs mass emerges as a geometric breathing mode and the neutrino mass scale via cubic geometric suppression. This work suggests that the 19 free parameters of the Standard Model are topological invariants of information processing.
\end{abstract}

\maketitle

The Standard Model's 19 free parameters lack theoretical derivation. We demonstrate these emerge from geometric constraints on a discrete, reflexive substrate where flux at time $t$ couples to its state at $t-\tau$.

Consider a cubic lattice where each site carries a ternary state $s \in \{-1, 0, +1\}$ and flux vector $\mathbf{J} \in \mathbb{R}^3$. The flux satisfies a Gauss constraint $\nabla \cdot \mathbf{J} = \rho_s$. The dynamics are governed by the Reflexive Lagrangian:
\begin{equation}
\mathcal{L} = \frac{1}{2}\left(\frac{\partial \mathbf{J}}{\partial t}\right)^2 - \Phi(\nabla \cdot \mathbf{J} - \rho) + \frac{\kappa}{2}\mathbf{J}(t) \cdot \mathbf{J}(t-\tau)
\end{equation}

\textbf{Theorem 1} (Fibonacci Spectral Quantization): The Euler-Lagrange equation for the zero-momentum mode is $\ddot{J} + \kappa(J(t) - J(t-\tau)) = 0$. With discrete time $\tau = 1$, this yields the recurrence $J(t) = J(t-1) + J(t-2)$, giving characteristic equation:
\begin{equation}
\lambda^2 - \lambda - 1 = 0
\end{equation}
whose roots are $\phi = (1+\sqrt{5})/2$ and $\psi = -1/\phi$. Eigenfrequencies scale as $\phi^n$, proving $n_{\text{eff}} = 13 = F_7$ and $55 = F_{10}$ are \textit{dynamical necessities}.

\textbf{Theorem 2} (Degrees of Freedom): On the minimal $2 \times 2 \times 2$ lattice:
\begin{equation}
N_{\text{DoF}} = 24 - 7 - 1 = 16
\end{equation}
counting flux components ($8 \times 3$) minus independent Gauss constraints minus gauge freedom.

\textbf{Theorem 3} (CM Selection): The 4-fold discrete symmetry of 4D spacetime ($i^4 = 1$) uniquely selects the Gaussian integers $\mathbb{Z}[i]$, hence the $j = 1728$ complex multiplication point. This is the lemniscate curve $y^2 = x^3 - x$, with modulus $k = 1/\sqrt{2}$ and period:
\begin{equation}
K(1/\sqrt{2}) = \frac{\Gamma(1/4)^2}{4\sqrt{\pi}}
\end{equation}

Combining these factors yields the geometric coupling:
\begin{equation}
G^* = \sqrt{2} \cdot \frac{\Gamma(1/4)^2}{2\pi} \approx 2.9587
\end{equation}

The gauge couplings satisfy the master quadratic:
\begin{equation}
x^2 - 16(G^*)^2 x + 16(G^*)^3 = 0
\end{equation}

The roots are $x_+ = 137.036$ and $x_- = 3.024$. We \textit{conjecture} $x_+ = \alpha^{-1}$ and $\lfloor x_- \rfloor = N_c = 3$.

The composition constant $K_{comp} = m_e/\pi$ represents the topological energy cost for 3D manifestation. Baryon masses follow $M_{\text{phys}} = M_{\text{geo}} - K_{comp}$, where geometric masses involve Fibonacci integers:
\begin{align}
M_p &= \left(\frac{13}{\alpha} + 55\right)m_e - \frac{m_e}{\pi} = 938.2724 \text{ MeV} \\
M_n &= M_p^{\text{geo}} + (\phi^2 - 12\alpha)m_e - \frac{m_e}{\pi} = 939.5654 \text{ MeV}
\end{align}

\textbf{The Delta Baryon} (first excited state) uses mode $17 = 13 + 4$ and correction $81 = 3^4$:
\begin{equation}
M_\Delta = \left(\frac{17}{\alpha} + 81\right)m_e - \frac{m_e}{\pi} = 1231.7 \text{ MeV}
\end{equation}
\textbf{Experimental: $1232 \pm 2$ MeV. Agreement: 0.03\%.}

The Higgs represents the scalar breathing mode of the $2^3$ unit cell with self-coupling $\lambda_H = 1/8$:
\begin{equation}
m_H = v/2 = 123.1 \text{ GeV} \quad \text{(Exp: 125.2 GeV, 1.7\%)}
\end{equation}

Neutrinos scale via cubic geometric suppression:
\begin{equation}
m_\nu = m_e \cdot \alpha^3 \approx 0.20 \text{ eV} \quad \text{(KATRIN: } < 0.8 \text{ eV)}
\end{equation}

\begin{table}[h]
\caption{Predictions vs. experiment}
\begin{ruledtabular}
\begin{tabular}{lccl}
Quantity & Predicted & Experimental & Status \\
\hline
$\alpha^{-1}$ & 137.0360 & 137.0360 & 1.3 ppm \\
$M_p$ (MeV) & 938.2724 & 938.2720 & 0.4 keV \\
$M_n$ (MeV) & 939.5654 & 939.5654 & $<$ 1 eV \\
$M_\Delta$ (MeV) & 1231.7 & 1232 & 0.03\% \\
$m_H$ (GeV) & 123.1 & 125.2 & 1.7\% \\
$m_\nu$ (eV) & 0.20 & $<$ 0.8 & Valid \\
\end{tabular}
\end{ruledtabular}
\end{table}

The derivation chain is: (1) Reflexive Lagrangian $\to$ Fibonacci modes [Theorem]; (2) Gauss constraint $\to$ 16 DoF [Theorem]; (3) 4D symmetry $\to$ $j=1728$ [Theorem]; (4) Master quadratic $\to$ $\alpha^{-1}$, $N_c$ [Conjecture]; (5) Composition constant $\to$ baryon masses [Conjecture].

The framework predicts exactly three fermion generations ($N_{\text{gen}} = \lfloor x_- \rfloor = 3$) and the Delta Baryon mass to 0.03\% without free parameters. Full derivations appear in the companion paper \cite{GSM2025}.

\begin{thebibliography}{3}
\bibitem{Morel2020} L. Morel \textit{et al.}, Nature \textbf{588}, 61 (2020).
\bibitem{PDG2022} R. L. Workman \textit{et al.} (PDG), PTEP \textbf{2022}, 083C01 (2022).
\bibitem{GSM2025} W. J. Steinmetz III, ``The Geometric Standard Model,'' (2026).
\end{thebibliography}

\end{document}
